\section{Agent Repository Controls and Tool Adapters}

Public agent tooling converges on one lesson: instructions matter, but deterministic control surfaces matter more. Codex supports project guidance through \texttt{AGENTS.md} \cite{openaiCodexAgentsMd2026}. The community \texttt{AGENTS.md} specification frames the file as a plain Markdown companion for coding agents \cite{agentsMdSpec2026}. Claude Code loads project memory through \texttt{CLAUDE.md} and related hierarchy \cite{claudeCodeMemoryDocs2026}. Cursor uses project rules with scope metadata \cite{cursorRulesDocs2026}. GitHub Copilot supports repository and path-specific custom instructions \cite{githubCopilotInstructions2026}. These mechanisms should not become divergent manuals. They should be generated or mirrored from one canonical standard.

The adapter problem is an authority problem, not a formatting problem. Jankurai treats prompt files as thin routers over machine-readable policy: \texttt{agent/owner-map.json} says who may own a path, \texttt{agent/test-map.json} says which lane proves it, \texttt{agent/generated-zones.toml} says which outputs are read-only unless regenerated from source, and \texttt{agent/standard-version.toml} binds the instruction surface to the standard and paper edition.

The source-of-authority hierarchy is:
\begin{center}
\small
\texttt{standard/version manifest -> owner/test/generated/security maps -> proof lanes -> adapter mirrors -> task prompt -> untrusted issue/PR/web/tool text}
\end{center}
For example, if \texttt{AGENTS.md} appears to permit generated edits but \texttt{agent/generated-zones.toml} denies hand edits, adapter verification fails and the generated-zone policy wins. Provider files may summarize policy, but they may not widen permissions.

\begin{table*}[t]
\caption{Agent adapter conformance checklist.}
\label{tab:adapter-checklist}
\centering
\small
\begin{tabularx}{\textwidth}{L{0.24\textwidth} Y}
\toprule
\textbf{Check} & \textbf{Required evidence} \\
\midrule
Root router points to canonical policy & \texttt{AGENTS.md} or equivalent names the standard brief, maps, generated zones, and validation lanes. \\
Provider files do not widen permissions & Claude, Cursor, Copilot, Codex, hooks, and MCP guidance agree with owner/test/security/generated maps. \\
Context packs label trust & Included material is labeled as trusted policy, repo code, generated artifact, proof evidence, docs, or untrusted input. \\
Adapters bind to policy digest/version & Mirrors include or can verify standard version, schema version, and policy digest. \\
Conflicts fail closed & Contradictions select the narrower authority and produce an adapter drift finding. \\
\bottomrule
\end{tabularx}
\end{table*}

\begin{table*}[t]
\caption{Agent tool adapter matrix.}
\label{tab:tool-adapters}
\centering
\small
\begin{tabularx}{\textwidth}{L{0.16\textwidth} L{0.21\textwidth} Y Y}
\toprule
\textbf{Tool family} & \textbf{Native surface} & \textbf{Jankurai adapter} & \textbf{Non-negotiable rule} \\
\midrule
Codex & \texttt{AGENTS.md}, local scoped instructions. & Root router plus local ownership files and mapped proof commands. & Do not duplicate durable policy outside the standard. \\
Claude Code & \texttt{CLAUDE.md}, project memory, local rules. & Mirror root router, version, commands, and generated zones. & Memory may summarize policy but not grant broader permissions. \\
Cursor & \texttt{.cursor/rules} with scope metadata. & Generate scoped rules from owner map and test map. & Glob rules must agree with owner-map paths. \\
GitHub Copilot & Repository and path-specific custom instructions. & Short repo-wide mirror plus path instructions for high-risk owners. & No private architecture fork in IDE guidance. \\
MCP and hooks & Tool schemas, hooks, and local automation. & Pin source, scope permissions, and separate untrusted data from trusted policy. & Tool output cannot rewrite authority. \\
Terminal agents & CLI prompt files, repo maps, shell wrappers. & Root router, filtered command output, raw-output escape hatch. & Tool convenience must not bypass CI gates. \\
\bottomrule
\end{tabularx}
\end{table*}

Token minimization is not a prompting trick; it is repository design. Large root instructions, duplicate architecture documents, stale docs, broad rereads, and noisy command output waste context and invite contradiction. ContextBench and SWE-Effi support this direction by treating useful context and resource constraints as first-class bottlenecks \cite{contextBench2026,sweEffi2025}. Context-pack metrics should include token count, stale docs excluded, proof-selection precision, route confidence, false-exclusion rate, and first-pass repair success.

The context-pack contract is budgeted and trust-labeled. A pack carries \texttt{--max-tokens}, an estimated token count, included and excluded file lists, source-trust labels, changed paths, selected owner routes, selected proof lanes, and evidence pointers. The result is not a larger prompt. It is a smaller grant of authority backed by files the repo can audit.
